\section{Conclusion}
\label{sec:conclusion}

Ratio is a bet that the mechanical core of accounting is a conservation
law and should be treated like one: stated as mathematics, proved once,
and relied on forever. We modeled accounting state as integer vectors
over a basis of conserved dimensions, defined validity as conservation
($\rvec = \rzero$), and defined the ledger as the monoidal fold of an
append-only log. The central guarantee --- that a ledger assembled from
balanced transactions is itself balanced (\cref{thm:conservation}) --- is
machine-checked in Lean~4, and classical double-entry falls out as the
one-dimensional instance (\cref{thm:double-entry}).

Three architectural commitments turn that theorem into a system.
\emph{Non-expressiveness} keeps the kernel small enough to prove and
cheap enough to run --- branch-minimal, fixed-width, integer-only, with a
cost dominated by memory bandwidth (\cref{sec:kernel}).
\emph{Emission} closes the proof-to-artifact gap: the running Rust is
generated from the proven Lean, so the deployed kernel inherits the proof
instead of approximating it (\cref{sec:emission}).
\emph{Externalization} relocates all business semantics to a
content-addressed control plane and all external data to typed facts with
explicit provenance, fencing learned components out of the trusted core
(\cref{sec:control}).

The result synthesizes ideas from programming languages, databases, and
computer architecture into one substrate: high assurance and high
throughput \emph{simultaneously}, with auditability and deterministic
replay as structural properties rather than bolted-on features. Treating
accounting as deterministic conservation over immutable inputs is what
makes all three compatible at once. The same recipe --- a proven,
non-expressive kernel; emission from proof to artifact; and expressiveness
externalized to a versioned, content-addressed plane --- applies to any
domain whose core is a conservation law wearing a costume of business
rules.
